% Original mathematical examples, written for this template demonstration.
% No excerpt from a particular textbook is reproduced here.
\title{Partial Differential\\Equations}
\subtitle{澄明讲义\enspace /\enspace A bilingual reading companion}
\author{Lumenbook · 学习笔记模板}
\date{9 November 2026}
\begin{document}
\frontmatter
\maketitle
\tableofcontents
\mainmatter

\chapter{Diffusion and energy}
\section{A model problem}
The heat equation connects a local differential law with a global
energy principle. We begin with a one-dimensional rod of length $L>0$,
with both endpoints held at zero temperature.\mn{\marginheading{Reading lens / 阅读线索}
先圈出空间区域、时间区间、初值与边界条件，再看作者使用的解的正则性。}
\begin{equation}\label{eq:heat}
  \begin{cases}
    u_t-u_{xx}=0, & 0<x<L,\quad 0<t<T,\\
    u(0,t)=u(L,t)=0, & 0\leq t\leq T,\\
    u(x,0)=u_0(x), & 0\leq x\leq L.
  \end{cases}
\end{equation}

\begin{definition}[Classical solution]\label{def:classical}
For this example, a classical solution has continuous derivatives
$u_t$, $u_x$ and $u_{xx}$ on $[0,L]\times[0,T]$ and satisfies
\eqref{eq:heat} pointwise, with one-sided derivatives at the endpoints.
\end{definition}

\begin{marginfigure}
  \begin{tikzpicture}[x=0.42cm,y=0.65cm]
    \draw[->,LumenMuted] (0,0)--(9.5,0) node[right] {$x$};
    \draw[->,LumenMuted] (0,0)--(0,2.7) node[above] {$u$};
    \draw[LumenAccent,thick,domain=0:9,samples=50]
      plot (\x,{2.1*sin(20*\x)});
    \draw[LumenBlue,dashed,domain=0:9,samples=50]
      plot (\x,{1.0*sin(20*\x)});
    \node[below] at (9,0) {$L$};
  \end{tikzpicture}
  \caption{A decaying sine mode. 振幅随时间衰减。}\label{fig:mode}
\end{marginfigure}
\begin{theorem}[Energy dissipation]\label{thm:energy}
For a classical solution of \eqref{eq:heat}, the energy
$E(t)=\frac12\int_0^L |u(x,t)|^2\dd x$ satisfies
\[
  E'(t)=-\int_0^L |u_x(x,t)|^2\dd x\leq 0.
\]
\end{theorem}
\begin{proof}
Multiply the equation by $u$ and integrate in space. Integration by
parts gives
\[
  E'(t)=\int_0^L u u_t\dd x
  =\bigl[u u_x\bigr]_0^L-\int_0^L |u_x|^2\dd x.
\]
The boundary term vanishes because $u=0$ at both endpoints.
\end{proof}
\begin{translationbox}
\textbf{中文理解。}先把方程乘以解本身，再做分部积分。
边界条件消去端点项，扩散项因此变成非正的能量变化率。
“能量不增加”来自这个符号，而不是温度处处单调下降。
\end{translationbox}

\newpage
\section{From derivatives to test functions}
For Poisson's equation $-\Delta u=f$ on a bounded domain
$\Omega\subset\R^d$, the weak formulation moves one derivative from
$u$ onto a test function.\sn{\label{sn:weak}Weak 不表示结论较弱；这里首先改变的是“满足方程”的表达方式。}
Under suitable smoothness and zero boundary conditions,
\begin{equation}\label{eq:weak}
  \int_\Omega \nabla u\cdot\nabla v\dd x
  =\int_\Omega fv\dd x.
\end{equation}

\begin{margintable}
  \caption{A compact glossary / 术语表}\label{tab:glossary}
  \begin{smalltable}
    \begin{tabular}{@{}ll@{}}
      \toprule
      English & 中文\\\midrule
      test function & 测试函数\\
      boundary trace & 边界迹\\
      coercivity & 强制性\\
      weak solution & 弱解\\\bottomrule
    \end{tabular}
  \end{smalltable}
\end{margintable}
\begin{definition}[Weak Dirichlet solution]\label{def:weak}
Let $\Omega$ be a bounded Lipschitz domain and $f\in L^2(\Omega)$.
A function $u\in H_0^1(\Omega)$ is a weak solution if
\eqref{eq:weak} holds for every $v\in H_0^1(\Omega)$.
\end{definition}

\begin{readingbox}{What to record beside a theorem}
\begin{enumerate}
  \item \textbf{Objects.} Which function space contains the unknown?
  \item \textbf{Hypotheses.} Where are boundedness and boundary regularity used?
  \item \textbf{Mechanism.} Is the argument based on energy, compactness or duality?
  \item \textbf{Scope.} Does the conclusion give existence, uniqueness or regularity?
\end{enumerate}
\end{readingbox}

\begin{margincode}[language=Python]
# 能量的离散近似
import numpy as np
E = 0.5 * dx * np.sum(u**2)
\end{margincode}
For numerical intuition, approximate the energy by a quadrature rule.
This computation illustrates a quantity; a discrete scheme needs its
own stability argument. A plot alone does not prove the energy identity.

\begin{example}[A single sine mode]
With $k=\pi/L$, the function
\[
  u(x,t)=e^{-k^2t}\sin(kx)
\]
solves \eqref{eq:heat} with $u_0(x)=\sin(kx)$. Substitution verifies
both the equation and the boundary values. Figure~\ref{fig:mode}
illustrates two amplitudes of this mode.
\end{example}
\begin{exercise}
For this explicit solution, compute $E(t)$ and check
Theorem~\ref{thm:energy} directly. Then explain why the sign in the
exponential must be negative.
\end{exercise}

\chapter{A notebook that can grow}
\section{Read, translate, reconstruct}
英文正文保留原有的数学叙述节奏，中文边注记录术语、直觉和局部疑问。
遇到一个重要结论，可以先整理为定义或定理，再用底色框写出自己的理解。
\mn{\textbf{本页验证}\\中文正文、中文边注、彩色框，以及带编号的交叉引用。
默认模式中，本页和前一页的边注都在右侧。}

\begin{notebox}{Three passes / 三遍阅读}
\textbf{第一遍：定位。}写出研究对象与结论，暂时跳过技术细节。\par
\textbf{第二遍：重建。}合上书，补全每一步推理需要的条件。\par
\textbf{第三遍：检验。}用一个简单模型和一个边界情形检查自己的理解。
\end{notebox}

\begin{proposition}[A useful reading distinction]\label{prop:reading}
A formal calculation suggests a proof strategy. To turn it into a
proof, justify every derivative, integral and limiting operation under
the hypotheses actually stated.
\end{proposition}
\begin{remark}
这里的命题用于演示排版。学习笔记应区分原书结论、自己补写的证明和
尚待确认的猜想，并在实际引用时注明原书版本、章节和页码。
\end{remark}

\begin{readingbox}{Cross-references that stay useful}
The class uses ordinary \LaTeX{} labels: see
Definition~\ref{def:classical}, Theorem~\ref{thm:energy},
equation~\eqref{eq:weak} and Table~\ref{tab:glossary}.
The numbered sidenote on page~\pageref{sn:weak} has its own label.
\end{readingbox}

\subsection{Rich material in the margin}
\begin{marginblock}
  \begin{tcolorbox}[lumen panel,breakable=false,title={小结 / Recap},
    left=5pt,right=5pt,top=5pt,bottom=5pt,before skip=0pt,after skip=0pt]
    边栏也能放小型底色框。\\[3pt]
    $\displaystyle E(t)\leq E(0)$\\[3pt]
    Keep one idea per note.
  \end{tcolorbox}
\end{marginblock}
Figures and tables share the normal counters. Margin blocks use their
own local line width, so \verb|width=\linewidth| adapts an image to
the sidebar. A small table can be wrapped in \verb|smalltable| to shrink
it only when its natural width exceeds the available space.

\begin{translationbox}
边栏用于补充主线。较长的证明、复杂表格和大段代码适合留在正文；
整页宽的公式或比较表可以使用后面的 \texttt{widepage} 环境。
\end{translationbox}

\newpage
\section{A small, predictable interface}
Keep the class, style files and font folder beside your main
\texttt{.tex} file. To start writing, use the supplied \texttt{main.tex}.
\mn{\textbf{偶数页检查}\\若以 \texttt{print} 模式编译，本页的这条边注应位于左侧。
正文版心同时向右移动，为装订后的外侧边栏留出空间。}

\begin{codeblock}[language=LaTeX,title={A minimal notebook}]
\documentclass{lumenbook} % 默认边注始终在右侧
% 双面打印时，改为：
% \documentclass[print,bindingoffset=4mm]{lumenbook}
\title{My PDE notebook}
\author{Your name}
\begin{document}
\maketitle
\chapter{The heat equation}
English text.\mn{这里写中文边注。}
\end{document}
\end{codeblock}

\subsection{Putting code in the sidebar}
Use \verb|margincode| for code in a margin. A language label and line
numbers make short examples easy to discuss. Chinese comments and
automatic wrapping are supported. File imports infer the language from
the extension, for example \verb|\inputcode{solver.py}|.
\begin{margincode}[language=Python,title={Heat step}]
# 显式差分：仅作代码排版示例
r = dt / dx**2
u[1:-1] += r * (
    u[2:] - 2*u[1:-1] + u[:-2]
)
\end{margincode}

\begin{notebox}{Placement controls / 边注位置}
Without an offset, notes use the standard margin-placement mechanism.
It separates adjacent notes.\par
\verb|\mn[6mm]{...}| and \verb|\begin{marginfigure}[6mm]| pin a note
with a vertical offset instead. Positive values move it down.
Pinned notes need a visual check because they may overlap.
\end{notebox}

\begin{readingbox}{Choose the appearance}
Choose jade, indigo or grayscale with the theme option. The language
option switches theorem and figure names to Chinese. Set the boxes
option to false for plain theorems. Chinese text is available in every mode.
\end{readingbox}

\begin{widepage}
\section*{A full-width comparison / 整页宽的比较表}
This page deliberately borrows the sidebar area. The following page
restores the original margin geometry automatically.
\begin{readingbox}{Keep the question, hypotheses and mechanism together}
Use a wide page for content that benefits from horizontal space.
No margin note is placed on this page.
\end{readingbox}
\begin{center}
\renewcommand{\arraystretch}{1.7}
\begin{tabularx}{\linewidth}{@{}l*{3}{>{\RaggedRight\arraybackslash}X}@{}}
\toprule
Viewpoint & Main object & What to check & 中文提示\\\midrule
Classical & Pointwise derivatives & Differentiability and boundary values
  & 每一步微分运算都有意义。\\
Weak & Integral identities & Test space and trace conditions
  & 导数转移给测试函数。\\
Energy & Norms and quadratic quantities & Signs and boundary terms
  & 估计往往比显式解更重要。\\
Numerical & A finite set of unknowns & Consistency, stability and convergence
  & 图像直觉与数学证明分开记录。\\\bottomrule
\end{tabularx}
\end{center}
\vspace{12mm}
\begin{equation*}
  \underbrace{\frac12\frac{\dd}{\dd t}\norm{u(t)}_{L^2(\Omega)}^2}
    _{\text{change in energy}}
  +\underbrace{\norm{\nabla u(t)}_{L^2(\Omega)}^2}
    _{\text{dissipation}}
  =\underbrace{0}_{\text{homogeneous heat equation}}.
\end{equation*}
\end{widepage}

\section{Back to the margin layout}
The margin width, headers and page-side behavior have been restored.
\mn{\textbf{恢复检查 / Restored layout}\\这是整页宽环境之后的第一条边注。
打印模式下本页是偶数页，边注应回到左侧；默认模式下仍在右侧。}
\begin{readingbox}{Ready for your own textbook}
Record the edition and page range before you begin. Then replace this
demonstration with your own chapter files, keeping English statements,
Chinese explanations and open questions close to one another.
\end{readingbox}
\begin{definition}[A useful note]
A useful note identifies one idea, states the assumptions it needs and
points back to the place where the idea is used.
\end{definition}
\end{document}
